%-------------------------
% Resume in Latex
% Author : Jose Vega
% License : MIT
%------------------------

\documentclass[letterpaper,11pt]{article}

\usepackage{latexsym}
\usepackage[empty]{fullpage}
\usepackage{titlesec}
\usepackage{marvosym}
\usepackage[usenames,dvipsnames]{color}
\usepackage{verbatim}
\usepackage{enumitem}
\usepackage[hidelinks]{hyperref}
\usepackage{fancyhdr}
\usepackage[english]{babel}
\usepackage{tabularx}

\pagestyle{fancy}
\fancyhf{} % clear all header and footer fields
\fancyfoot{}
\renewcommand{\headrulewidth}{0pt}
\renewcommand{\footrulewidth}{0pt}
\setlength{\footskip}{4.08003pt}

% Adjust margins
\addtolength{\oddsidemargin}{-0.5in}
\addtolength{\evensidemargin}{-0.5in}
\addtolength{\textwidth}{1in}
\addtolength{\topmargin}{-.5in}
\addtolength{\textheight}{1.0in}

\urlstyle{same}

\raggedbottom
\raggedright
\setlength{\tabcolsep}{0in}

% Sections formatting
\titleformat{\section}{
  \vspace{-4pt}\scshape\raggedright\large
}{}{0em}{}[\color{black}\titlerule \vspace{-5pt}]

%-------------------------
% Custom commands
\newcommand{\resumeItem}[2]{
  \item\small{
    \textbf{#1}{: #2 \vspace{-2pt}}
  }
}

\newcommand{\resumeSubheading}[4]{
  \vspace{-1pt}\item
    \begin{tabular*}{0.97\textwidth}[t]{l@{\extracolsep{\fill}}r}
      \textbf{#1} & #2 \\
      \textit{\small#3} & \textit{\small #4} \\
    \end{tabular*}\vspace{-5pt}
}

\newcommand{\resumeSubSubheading}[2]{
    \begin{tabular*}{0.97\textwidth}{l@{\extracolsep{\fill}}r}
      \textit{\small#1} & \textit{\small #2} \\
    \end{tabular*}\vspace{-5pt}
}

\newcommand{\resumeSubItem}[2]{\resumeItem{#1}{#2}\vspace{-4pt}}

\renewcommand{\labelitemii}{$\circ$}

\newcommand{\resumeSubHeadingListStart}{\begin{itemize}[leftmargin=*]}
\newcommand{\resumeSubHeadingListEnd}{\end{itemize}}
\newcommand{\resumeItemListStart}{\begin{itemize}}
\newcommand{\resumeItemListEnd}{\end{itemize}\vspace{-5pt}}

%-------------------------------------------
%%%%%%  Resume STARTS HERE  %%%%%%%%%%%%%%%%%%%%%%%%%%%%


\begin{document}
%----------HEADING-----------------
\begin{tabular*}{\textwidth}{l@{\extracolsep{\fill}}r}
  \textbf{\href{https://linkedin.com/in/josevegar}{\Large José E. Vega Raymondi}} & Email : \href{mailto:jvegar@uni.pe}{jvegar@uni.pe}\\
  \href{https://github.com/jvegar}{https://github.com/jvegar} & Mobile : +51-991139451 \\
\end{tabular*}

%-----------EDUCATION-----------------
\section{Educación}
  \resumeSubHeadingListStart
    \resumeSubheading
      {Universidad del Pacífico}{Lima, Perú}
      {Magister en Finanzas - Gestión de Portafolios de Inversión}{Dic. 2014 -- Dic. 2016}
    \resumeSubheading
      {Universidad Nacional de Ingeniería}{Lima, Perú}
      {Bachiller en Ingeniería de Sistemas}{Ene. 2003 -- Dic. 2009}
  \resumeSubHeadingListEnd


%-----------EXPERIENCE-----------------
\section{Experiencia}
  \resumeSubHeadingListStart

\resumeSubheading
  {BairesDev}{California, EEUU}
  {Software Engineer}{Set 2021 - Actualidad (Remoto)}
  \resumeItemListStart
    \resumeItem{Desarrollo Backend}
      {Modelamiento de BD (Postgres), Desarrollo de APIs (Node, JS, TypeScript, GraphQL, REST), Frameworks (NestJS, Express, Apollo Server, Hapi)}   
    \resumeItem{Desarrollo Front-End}
      {Prototipado UI (Figma), Desarrollo de componentes UI (Angular, React, HTML, CSS, JS, TS).}                  
    \resumeItem{DevOps}
      {Uso de Tablero Agile (Jira), Manejo de repositorios GIT, uso de contenedores (Docker) y servicios de la nube (AWS/Azure).}

  \resumeItemListEnd

  \resumeSubheading
  {NTT Data}{Lima, Perú}
  {Centers Lead Developer}{Oct 2020 - Ago 2021}
  \resumeItemListStart
    \resumeItem{Liderazgo}
      {A cargo de equipo de Microsoft para desarrollo de evolutivos y correctivos en los distintos clientes. Uso de metodología agil (planning, dailys, sprint reviews)}   
    \resumeItem{Arquitectura de datos}
      {Modelamiento de BD (Azure SQL, Postgres, Dataverse), diseño de tableros de control (PowerBI).}              
    \resumeItem{Desarrollo de Aplicaciones}
      {Desarrollo de componentes servidor (SharePoint 2013/2016), desarrollo de componentes FrontEnd (JS, React), desarrollo de componentes servidor (NodeJS, .NET Core).}
   
  \resumeItemListEnd

    \resumeSubheading
      {Certero}{Lima, Perú}
      {Arquitecto de Soluciones}{Feb 2017 - Set 2020}
      \resumeItemListStart
        \resumeItem{Arquitectura de datos}
          {Modelamiento de BD (SQL Server, Postgres), construcción de ETLs (SSIS, Azure Data Factory), diseño de tableros de control (PowerBI).}
        \resumeItem{Desarrollo de Aplicaciones Web Full Stack}
          {Desarrollo de componentes servidor (SharePoint 2013/2016/2019), desarrollo de componentes Front-End (HTML,CSS,JavaScript, Node), aplicaciones MVC .NET.}          
        \resumeItem{Arquitectura de Aplicaciones}
          {Diseño de las capas de aplicaciones y negocio. Mapeo de procesos (BMPN), Consumo y construcción de REST APIs (.NET Core, NodeJS).}
 
      \resumeItemListEnd
      
% --------Multiple Positions Heading------------
%    \resumeSubSubheading
%     {Software Engineer I}{Oct 2014 - Sep 2016}
%     \resumeItemListStart
%        \resumeItem{Apache Beam}
%          {Apache Beam is a unified model for defining both batch and streaming data-parallel processing pipelines}
%     \resumeItemListEnd
%    \resumeSubHeadingListEnd
%-------------------------------------------

    \resumeSubheading
      {Evol TS Net}{ Lima, Perú}
      {Consultor Desarrollador}{Ene 2014 - Ene 2017}
      \resumeItemListStart
        \resumeItem{Desarrollo Backend}
          {Construcción de componentes de lado servidor (webparts, event receivers, timer jobs). Desarrollo de Aplicaciones MVC y servicios WCF de integración con SAP.}
        \resumeItem{Base de Datos}
          {Modelamiento de base de datis, Construcción de componentes de BD (T-SQL), Optimización de queries (índices, tablas temporales).}        
        \resumeItem{Desarrollo Front-End}
          {Rediseño de componentes Front-End (Páginas maestra) utilizando metodologías de UX, diseño responsivo (Bootstrap) y librerías JS (JQuery, Knockout).}
      \resumeItemListEnd

      \resumeSubheading
      {CGI}{Lima, Perú}
      {Analista Programador}{Dic 2011 - Dic 2013}
      \resumeItemListStart
        \resumeItem{Backend Development}
          {Optimización de servidor de aplicacione (IIS), construcción de componentes de lado servidor (webparts, scripts Powershell, servicios WCF).}
        \resumeItem{Base de datos}
          {Construcción de ETLs para integración entre base de datos y archivos CSV (SQL Server).}
      \resumeItemListEnd

%    \resumeSubheading
%      {CGI}{Lima, Perú}
%      {Analista Programador}{Dic 2011, Dic 2013}
%      \resumeItemListStart
%        \resumeItem{SharePoint Server}
%          {Configuración de granja SharePoint, instalación de nuevos servidores. Construcción de funcionalidades.}
%        \resumeItem{SQL Server Integration Services}
%          {Construcción de ETL para la integración de bases de datos SQL y archivos CSV.}
%      \resumeItemListEnd

%    \resumeSubheading
%      {Banco de la Nación}{Lima, Perú}
%      {Analista de Infraestructura de TI}{Jun 2009 - Jul 2011}
%      \resumeItemListStart
%        \resumeItem{IBM Portal Server}
%          {Construcción de Portlets utilizando Java, y construcción de Web Services SOAP.}
%        \resumeItem{IBM Business Process Manager}
%          {Implementación de Pruebas de Concepto de Flujo de Proceso BPM.}
%      \resumeItemListEnd

  \resumeSubHeadingListEnd


%-----------PROJECTS-----------------
\section{Proyectos}
  \resumeSubHeadingListStart
    \resumeSubItem{Implementación del Case Manager - Lucinity}
      {Construcción de APIs GraphQL usando TypeScript, Express, Apollo Server con NodeJS y de componentes React para la aplicación Case Manager.}
      \resumeSubItem{Migración de Portal Learning - Instructure}
      {Se migró el backend construido en Node v5 a v16. Se actualizaron las migraciones SQL Postgres de v9 a v12. Se reemplazaron las librerias JS obsoletas. Refactoring general para uso de features ES6 (let/const, promises, async/await)}
    \resumeSubItem{Implementación de Portal Institucional - Fonafe}
      {Desarrollo de CMS a medidad utilizando .NET Core, REST APIs, y MVC para el Backend, y Angulas JS y Bootstrap para Front-End.}
  \resumeSubHeadingListEnd

%
%--------PROGRAMMING SKILLS------------
\section{Habilidades}
  \resumeSubHeadingListStart
  \resumeSubItem{Languages}{: TypeScript, Javascript, CSharp, Java, Python }          
  \resumeSubItem{Frameworks}{: Express, NestJS, Apollo GraphQL, Hapi, Net Core, Spring Boot}      
  \resumeSubItem{Methodologies}{: Waterfall, Agile, Kanban, DevOps}      
  \resumeSubItem{Idioms}{: Espanol - Nativo / Inglés - Avanzado C1} 
  \resumeSubHeadingListEnd


%-------------------------------------------
\end{document}